% paper/sections/11c_dccd_bootstrap.tex
% ═════════════════════════════════════════════════════════════════════════════
% §11 統合更新手順周辺（DCCD 後置フィルタ設計は §4c に集約済；
%        Bootstrap は付録 H に集約済；本ファイルは初期整合化とタイムステップ制御のみ保持）
% ═════════════════════════════════════════════════════════════════════════════

\noindent
本節は時間発展を開始する前の幾何・圧力・格子計量係数の整合化と，
その後の $\Delta t$ 制御だけを扱う．
DCCD 後置フィルタは §\ref{sec:dccd_derivation} の節点中心 CCD 外殻設計であり，
§11 の標準更新手順では CLS 移流，圧力面上補正，または二相圧力ジャンプ系の
主演算子として再導入しない．
本章の演算子割当は §\ref{sec:scheme_per_variable} に従い，
二相圧力閉包は分相 PPE + HFE + DC + 毛管補正量で閉じる．

% ═════════════════════════════════════════════════════════════════════════════
\subsection{初期整合化：ブートストラップと時刻0の界面面幾何}
\label{sec:bootstrap}
% ═════════════════════════════════════════════════════════════════════════════

時間発展を開始する前に，格子--界面の循環参照を解消する
6 段階ブートストラップ（一様格子仮 SDF $\to$ 非一様格子生成 $\to$ CLS 初期化 $\to$
CCD 演算子定義 $\to$ 非一様格子上 SDF 再初期化 $\to$ Young--Laplace 初期圧力）で
一度だけ整合させる．詳細な手順は付録~\ref{app:bootstrap} を参照．
この初期整合化により，第1段に入る時点で
$\psi$，$\phi$，物性値，格子計量係数，初期圧力が同じ界面幾何を参照する．
格子の時間更新モード（固定格子モード，再生成モード）は §~\ref{sec:grid_ale} に示す．

% ═════════════════════════════════════════════════════════════════════════════
\subsection{タイムステップ制御}
\label{sec:algo_timestep}
% ═════════════════════════════════════════════════════════════════════════════

各タイムステップの開始時に，以下の2制約のうち最も厳しい条件で $\Delta t$ を決定する
（詳細な導出は §~\ref{sec:stability} 参照）：

\begin{equation}
  \Delta t = \min\!\left(
    \underbrace{
      C_{\mathrm{CFL}}
      \left[\max\!\left(\frac{|u|_{\max}}{\Delta x}
                       +\frac{|v|_{\max}}{\Delta y}\right)\right]^{-1}
    }_{\Delta t_{\mathrm{adv}}\text{：対流 CFL（式~\eqref{eq:dt_adv}）}},\;
    \underbrace{\sqrt{\frac{(\rho_l+\rho_g)\,h_{\min}^3}{2\pi\sigma}}}_{\Delta t_\sigma\text{：毛管波制約~\cite{DennerVanWachem2015}（式~\eqref{eq:dt_sigma}）}}
  \right)
  \label{eq:dt_full_algo}
\end{equation}

ここで $C_{\mathrm{CFL}} \le 0.5$ は安全率であり，本稿では 0.45 を用いる．
$h_{\min} = \min(\Delta x, \Delta y)$ である．
表面張力なし（$\sigma=0$）の計算では毛管波制約を除外し，対流 CFL を用いる．
粘性項は標準経路では陰的 BDF2 + 欠陥補正の全応力 Helmholtz 系で扱うため，
粘性 CFL は明示的安定制約としては現れない．
ただし大粘性比（$\mu_l/\mu_g \gg 1$）の気液系で界面近傍のクロス応力を
低次に外へ出す簡略化を行う場合は，実質的な粘性制約が復活する
（§~\ref{sec:stability} 参照）．

\noindent\textbf{3 項 $\min$ 版（陽的粘性領域）：}
式~\eqref{eq:dt_full_algo} の 2 項制約は，粘性剛性を陰的 Helmholtz 系が吸収する標準領域での縮約表現である．
陽的粘性領域では $\Delta t_\text{visc} = C_\text{visc}\rho h_{\min}^2/(2\mu N_\text{dim})$（粘性 CFL）も加え，
$\Delta t = \min(\Delta t_\text{adv}, \Delta t_\text{visc}, \Delta t_\sigma)$ の 3 項 $\min$ となる．
物理波速度 $\Delta t_\text{wave}$ は非圧縮解法のため両式とも除外される．

\noindent\textbf{非一様格子での局所版．}
非一様格子では $h_{\min}$ による式~\eqref{eq:dt_full_algo} は安全側の縮約である．
局所 $\varepsilon(\bx)$，高粘性比，格子再生成を併用する場合は，
\[
  \Delta t =
  \min\!\left(
    \Delta t_{\mathrm{adv,local}},
    \Delta t_{\sigma,\min},
    \Delta t_{\nu,\mathrm{interface}},
    \Delta t_{\mathrm{grid}}
  \right)
\]
を別の制約評価として併用する．
$\Delta t_{\sigma,\min}$ は界面帯の最小セル幅で評価し，
$\Delta t_{\nu,\mathrm{interface}}$ は $\mu(\bx),\rho(\bx)$ の界面近傍最大比を用いる．
再生成モードではさらに格子移動量が §~\ref{sec:grid_ale} の $0.5h$ 条件を越えないことを
$\Delta t_{\mathrm{grid}}$ で確認する．
ただしこの制約は，式~\eqref{eq:ale_grid_gcl} の離散幾何保存則，
式~\eqref{eq:ale_conservative_update} の保存更新，および
式~\eqref{eq:regrid_face_geometry} の界面面データ再構築が同時に成立する場合にのみ有効である．
これらが揃わない追随格子は本節の 1 タイムステップ標準経路へ含めない．

\bigskip
以上で標準経路の 1 タイムステップ更新は，
界面輸送，幾何更新，運動量予測，圧力投影，速度補正を
同一の界面面幾何上で閉じる形になる．
次章では，この統合更新手順を支える各構成要素が
設計通りの精度を達成するかを系統的に検証する．
